\documentclass[11pt]{article}
\usepackage[margin=1in]{geometry}
\usepackage[T1]{fontenc}
\usepackage[utf8]{inputenc}
\usepackage{hyperref}
\usepackage{xcolor}
\usepackage{enumitem}
\usepackage{titlesec}
\usepackage{fancyvrb}
\hypersetup{colorlinks=true, linkcolor=blue, urlcolor=blue}

\titleformat{\section}{\large\bfseries}{}{0em}{}
\title{Summary Generative AI Usage}
\author{Eduardo Rodríguez Sánchez}
\date{\today}

\begin{document}
\maketitle

This document contains a record of the technical prompts and AI-assisted refinements used during the development of the Communities and Crime visualization project.

\section{Layout overflow in the visualization container}
\subsection*{Prompt}
The visualization container overflows downward in the window. My current JSX structure is:
\begin{Verbatim}[fontsize=\small]
<div className="visualization-container">
  <div className="controls">...</div>
  <div ref={containerRef} className="visualization-container" />
</div>
\end{Verbatim}

\subsection*{AI Answer}
The issue is that the outer div and the inner div both have the same class, and the inner mount target is growing without constraint. 
\begin{itemize}[leftmargin=*]
  \item Remove the duplicate class on the outer wrapper.
  \item Use a flex column wrapper with \texttt{height: 100\%}.
  \item Set the D3 target to \texttt{flex: 1; min-height: 0; overflow: hidden;}. The \texttt{min-height: 0} is critical to allow flex children to shrink below their content size.
\end{itemize}

\section{SVG stacking and label visibility}
\subsection*{Prompt}
I tried adding a \texttt{z-index} to the state labels, but they still appear behind the circles.
\begin{Verbatim}[fontsize=\small]
.state-label { z-index: 2000; }
\end{Verbatim}

\subsection*{AI Answer}
\texttt{z-index} does not work on SVG elements. SVG follows the painter's model, where the DOM order determines stacking. To ensure labels appear on top, you should render them in a separate group appended after all the circles are drawn.

\section{Adding background boxes to SVG labels}
\subsection*{Prompt}
I want to put a semi-transparent box behind each label to make it readable over the colored circles. How do I size the box dynamically?

\subsection*{AI Answer}
You can use the \texttt{getBBox()} method to measure the rendered text and then insert a \texttt{rect} behind it:
\begin{Verbatim}[fontsize=\small]
const bbox = text.node().getBBox();
const pad = { x: 8, y: 4 };

labelG.insert("rect", "text")
  .attr("x", bbox.x - pad.x)
  .attr("y", bbox.y - pad.y)
  .attr("width", bbox.width + pad.x * 2)
  .attr("height", bbox.height + pad.y * 2)
  .attr("fill", "rgba(26, 29, 39, 0.75)");
\end{Verbatim}
Ensure you use \texttt{.datum(d)} when creating the label groups so that subsequent interaction updates can access the bound data.

\section{Refactoring shared label rendering}
\subsection*{Prompt}
How can I refactor the label logic so that both the treemap and the circle pack layouts use the same implementation?

\subsection*{AI Answer}
Extract a shared \texttt{renderLabels(data)} method. The only differences are the coordinate systems. You can use a conditional check on \texttt{this.currentLayout} to set the positions:
\begin{Verbatim}[fontsize=\small]
const x = this.currentLayout === "treemap" ? d.x0 + (d.x1 - d.x0) / 2 : d.x;
const y = this.currentLayout === "treemap" ? d.y0 + (d.y1 - d.y0) / 2 : d.y;
\end{Verbatim}
This allows you to maintain one set of styling and transition logic for both layouts.

\section{Enriching the dataset with state codes}
\subsection*{Prompt}
I have the FIPS state codes, but I want to display the two-letter abbreviations (CA, NY, etc.). How should I add this to the Redux slice?

\subsection*{AI Answer}
Create a mapping object from state names to codes and apply it inside the \texttt{.map()} function where you process the incoming CSV data:
\begin{Verbatim}[fontsize=\small]
const stateCode = STATE_TO_CODE[stateName] || "??";
return { ...item, state: stateName, stateCode: stateCode };
\end{Verbatim}

\section{Grouping hierarchical data with custom properties}
\subsection*{Prompt}
When using \texttt{d3.groups()} to create the hierarchy, how do I attach the state code to the parent node?

\subsection*{AI Answer}
Since all items in a state group share the same code, you can grab it from the first item in the group array:
\begin{Verbatim}[fontsize=\small]
.groups(data, (d) => d.state)
  .map(([state, values]) => ({
    name: state,
    stateCode: values[0].stateCode,
    children: values.map((v) => ({ ...v, name: v.communityname })),
  }));
\end{Verbatim}

\section{Centralizing interaction state logic}
\subsection*{Prompt}
My \texttt{update()} method is getting messy with many \texttt{if} statements for hovered, selected, and highlighted states. How can I make it more consistent?

\subsection*{AI Answer}
Create a \texttt{getNodeState(d)} helper to centralize the priority logic. This allows you to compute the state once and reuse it across stroke, fill, and opacity attributes:
\begin{Verbatim}[fontsize=\small]
const getNodeState = (d) => {
  if (d.data.id === hoveredCommunity) return "hovered";
  if (selectedCommunities.includes(d.data.id)) return "selected";
  if (d.data.state === highlightedState) return "highlighted";
  return "default";
};
\end{Verbatim}

\section{Solving click event interception}
\subsection*{Prompt}
The click handler on the state nodes isn't working because the leaf nodes (communities) are on top. How do I fix this?

\subsection*{AI Answer}
Move the click handler to the parent \texttt{<g>} group instead of the shapes. Clicks on children will bubble up to the group. Then, use \texttt{e.stopPropagation()} on leaf-level hovers if you want to prevent them from interfering with state-level selections.

\section{useMemo vs useEffect for derived state}
\subsection*{Prompt}
Should I use \texttt{useMemo} or \texttt{useEffect} to derive the \texttt{stateCode} for the badge in my React component?

\subsection*{AI Answer}
\texttt{useMemo} is correct here. It derives a value from existing state synchronously during render. \texttt{useEffect} would trigger an unnecessary second render cycle and extra state management.

\section{Responsive sizing with CSS clamp()}
\subsection*{Prompt}
The floating badge doesn't adapt to different screen sizes. How can I make it responsive without multiple media queries?

\subsection*{AI Answer}
Use the \texttt{clamp()} function for fluid typography and padding. It sets a minimum, a preferred value (using viewport units), and a maximum:
\begin{Verbatim}[fontsize=\small]
.state-badge__code {
  font-size: clamp(14px, 3vw, 28px);
}
\end{Verbatim}

\section{Tooltip placement in screen quadrants}
\subsection*{Prompt}
How do I make the tooltip flip its position so it always stays within the screen bounds?

\subsection*{AI Answer}
Compare the cursor position against the window midpoints and adjust the offsets accordingly:
\begin{Verbatim}[fontsize=\small]
const isTopHalf = event.pageY < window.innerHeight / 2;
const isLeftHalf = event.pageX < window.innerWidth / 2;

this.tooltip
  .style("top", isTopHalf ? (pageY + 10) : (pageY - height - 10))
  .style("left", isLeftHalf ? (pageX + 10) : (pageX - width - 10));
\end{Verbatim}

\section{Scaling SVG elements from center}
\subsection*{Prompt}
When I try to scale a node on highlight using \texttt{transform: scale(1.15)}, it scales from the top-left origin. How do I scale from the center?

\subsection*{AI Answer}
SVG \texttt{scale()} always uses the coordinate system origin. You must translate to the center of the element, scale, and then translate back:
\begin{Verbatim}[fontsize=\small]
const cx = (d.x1 - d.x0) / 2;
const cy = (d.y1 - d.y0) / 2;
return `translate(${d.x0},${d.y0}) translate(${cx},${cy}) 
        scale(1.15) translate(${-cx},${-cy})`;
\end{Verbatim}

\end{document}
